\documentclass[11pt]{article}

\usepackage[margin=1in]{geometry}
\usepackage[T1]{fontenc}
\usepackage{lmodern}
\usepackage[colorlinks=true,urlcolor=blue]{hyperref}

\pagestyle{empty}
\setlength{\parindent}{0pt}
\setlength{\parskip}{0.8em}

\begin{document}

\begin{flushright}
Saveliy Baturin\\
Independent Researcher\\
\href{mailto:saveliy.xo@gmail.com}{saveliy.xo@gmail.com}\\
August 30, 2026
\end{flushright}

Dear Editors of the SIAM Journal on Mathematics of Data Science,

Please consider the manuscript ``Hierarchical Convexity-Run Decomposition of
Sampled Curves: Finite-Sample Recovery and Stability.'' The paper introduces a
nonlinear multiscale representation for sampled curves based on runs of discrete
curvature sign. Its main contributions are an exactly reconstructing nested
transform with a sparse linear-work implementation, a precise characterization
of the discontinuity of hard knot decisions together with stable companion
summaries, and finite-sample recovery guarantees for the hierarchy and for an
explicit alternating chord-lobe model.

I believe the manuscript is particularly well suited to SIMODS because its
central question lies in mathematical data science: how to construct, analyze,
and statistically certify an interpretable representation whose atoms are
localized convex and concave departures from chords. The theory is complemented
by simulations that test the recovery boundary and by a no-refit transfer study
on expert-labelled chromatographic curves. That application serves as a
downstream validation setting for a general sampled-curve representation, with
matched controls and negative boundary experiments reported explicitly.

Code, protocols, compact results, and reproducibility instructions are archived
on Zenodo at
\href{https://doi.org/10.5281/zenodo.22171977}{doi:10.5281/zenodo.22171977};
active development is available at
\href{https://github.com/Safreliy/hcrd}{github.com/Safreliy/hcrd}. The manuscript
is original and is not under consideration elsewhere.

Sincerely,

Saveliy Baturin\\
Independent Researcher

\end{document}
